\begin{table}[t]
\centering
\caption{처치 상태별 기술통계}
\label{tab:descriptives_ko}
\begin{tabular}{lrrrr}
\toprule
변수 & 통제 (D=0) & 처치 (D=1) & 차이 (T-C) & p-값 \\
\midrule
op\_cost\_ex\_rent & 24,355,776 (48,046,410) & 8,309,086 (34,766,403) & -16,046,690 & <0.001 \\
농업경영비 전체 (원) & 26,280,404 (49,518,866) & 8,447,763 (34,898,052) & -17,832,642 & <0.001 \\
농외소득 (원) & 7,267,203 (12,614,381) & 11,339,222 (13,170,071) & 4,072,019 & <0.001 \\
가계 소비지출 (원) & 23,425,658 (10,712,154) & 22,353,227 (10,920,864) & -1,072,431 & 0.053 \\
농업소득 (원) & 15,821,036 (33,898,987) & 3,657,815 (19,021,054) & -12,163,221 & <0.001 \\
rent\_cost & 1,924,629 (4,334,323) & 138,677 (599,465) & -1,785,952 & <0.001 \\
2018 기준 경지면적 (처치 정의용, ㎡) & 14,967 (18,004) & 2,729 (1,169) & -12,238 & <0.001 \\
자작면적 (㎡) & 6,496 (4,163) & 1,962 (1,363) & -4,533 & <0.001 \\
Imputed 공익직불금 (면적 기준, 원) & 0 (0) & 0 (0) & 0 &  \\
type\_fulltime (mode) & 1 & 1 & --- & --- \\
sex\_code (mode) & 1 & 1 & --- & --- \\
edu\_code (mode) & 11 & 11 & --- & --- \\
관측치 (기준연도) & 1,311 & 867 & --- & --- \\
\bottomrule
\end{tabular}
\\
\footnotesize\textit{주: FHES 전국 가중값 적용. 가중 평균 (괄호 안 표준편차). p-값은 가중 t-검정. 기준연도 = 2018 단면.}
\end{table}

